% !PRINTED 26
\weberpage{26}
%<<CONT
und die Funktionen $\Phi$, $\Psi$ erhalten die kanonische Form:
\weq{14}{\begin{aligned}
\Phi &= \Sigma\,\beta_1^2\alpha_2\alpha_3\alpha_4\,y_1^2,\\
\Psi &= \Sigma\,\alpha_1^2\beta_2\beta_3\beta_4\,y_1^2.
\end{aligned}}
Die Determinante des Systems (12), (14), als lineare Gleichungen für
$y_1^2$, $y_2^2$, $y_3^2$, $y_4^2$ betrachtet, ist
\weq{15}{\begin{vmatrix}
\alpha_1^2\beta_2\beta_3\beta_4, & \alpha_2^2\beta_1\beta_3\beta_4, & \alpha_3^2\beta_1\beta_2\beta_4, & \alpha_4^2\beta_1\beta_2\beta_3\\
\beta_1^2\alpha_2\alpha_3\alpha_4, & \beta_2^2\alpha_1\alpha_3\alpha_4, & \beta_3^2\alpha_1\alpha_2\alpha_4, & \beta_4^2\alpha_1\alpha_2\alpha_3\\
\alpha_1, & \alpha_2, & \alpha_3, & \alpha_4\\
\beta_1, & \beta_2, & \beta_3, & \beta_4
\end{vmatrix}.}
Sie verschwindet, wenn $\alpha_1:\alpha_2 = \beta_1:\beta_2$ wird, weil
dann die beiden ersten Kolonnen miteinander proportional werden, und
ebenso wenn $\alpha_i:\alpha_k = \beta_i:\beta_k$ wird. Es ist daher
(15) als ganze Funktion der $\alpha$, $\beta$ betrachtet, teilbar durch
\weq{16}{\alpha_i\beta_k - \alpha_k\beta_i = (\alpha_i\beta_k),}
und folglich auch durch das Produkt aller dieser Faktoren
\weq{17}{(\alpha_1\beta_2)(\alpha_1\beta_3)(\alpha_1\beta_4)(\alpha_2\beta_3)(\alpha_2\beta_4)(\alpha_3\beta_4) = \Delta.}
Aus der Vergleichung der Grade ergibt sich, daß sich die beiden
Funktionen (15), (17) nur durch einen Zahlenfaktor unterscheiden
können, und dieser ergibt sich aus der Annahme $\beta_1 = 0$,
$\alpha_2 = 0$, $\alpha_4 = 0$ gleich 1. Die Determinante (15) ist
also dem Produkt $\Delta$ gleich, und das Quadrat von $\Delta$ ist die
Diskriminante $D$ der Funktion $f(\xi,\eta)$, also gleich der
Funktion $D$, von der wir angenommen haben, daß sie von Null
verschieden sei. Es lassen sich also die $y_1^2$, $y_2^2$, $y_3^2$,
$y_4^2$ linear durch $\varphi$, $\psi$, $\Phi$, $\Psi$ darstellen.

\begin{center}
\rule{2cm}{0.4pt}
\end{center}

Die Kovarianten $\Phi$, $\Psi$ ändern sich, wenn $\varphi$ und $\psi$
nach (4) durch $\varphi'$, $\psi'$ ersetzt werden. Sie gehören also
nicht zu der Grundkurve, sondern zu dem Formenpaar $\varphi$, $\psi$.
Aber es gehen bei dieser Substitution $\Phi$ und $\Psi$ als lineare
Funktionen von $y_1^2$, $y_2^2$, $y_3^2$, $y_4^2$ in lineare
Funktionen von $\varphi$, $\psi$, $\Phi$, $\Psi$ über. Diese Ausdrücke
sind ziemlich kompliziert. Wir brauchen sie hier aber nur für die
Punkte der Grundkurve, also für $\varphi = 0$, $\psi = 0$, und unter
dieser Voraussetzung werden sie sehr einfach. Es ist da nämlich
\weq{18}{\begin{aligned}
\Phi' &= \Sigma\,(p\alpha_1+q\beta_1)^2(m\alpha_2+n\beta_2)(m\alpha_3+n\beta_3)(m\alpha_4+n\beta_4)\,y_1^2,\\
\Psi' &= \Sigma\,(m\alpha_1+n\beta_1)^2(p\alpha_2+q\beta_2)(p\alpha_3+q\beta_3)(p\alpha_4+q\beta_4)\,y_1^2,
\end{aligned}}
%CONT>>
